
\subsubsection{Key Findings}

\textbf{mypy augmentation provides substantial but mechanism-agnostic benefit.} The primary finding is a +8.94pp pass@1 improvement from adding mypy to execution-based repair loops on HumanEval+ (GPT-4o-mini, k=5, three seeds consistent). The type-specificity hypothesis --- that mypy helps specifically for type-error problems --- is not supported: both conditions achieve identical 90.9\% repair rates on type-error problems (h-m2, n=22, differential=0.000). The most parsimonious interpretation --- pending ablation confirmation (see L2) --- is general diagnostic context enrichment: mypy output adds structured text to the repair prompt that may improve GPT-4o-mini's repair quality across all problem types, not specifically where mypy detects a type error. If confirmed, this interpretation would suggest that other structured diagnostic channels (linters, semantic analyzers) might provide similar benefits independent of whether they target the specific failure mode.

\textbf{Benchmark-level signal divergence constrains applicability.} The 70\% versus 0\% mypy error rate between HumanEval+ and MBPP+ is a new empirical finding with practical implications. HumanEval+ failures are dominated by undefined-name errors (\texttt{name-defined} category); MBPP+ failures appear to be pure algorithmic correctness errors invisible to mypy. A practitioner should profile the mypy error rate in failing solutions before adding mypy feedback --- if it is below 10\%, mypy will not activate frequently enough to drive improvement.

\textbf{Z3 counterexample feedback: sound architecture, insufficient activation.} The Z3 null result ($-0.9pp)$ is informative rather than merely negative. The spec generation pipeline achieves 91.1\% validity, confirming that LLM-generated Z3 specs are mostly correct on arithmetic HumanEval+ problems. The bottleneck is CE activation (16.1\%): most problems lack a Z3 counterexample to provide, making Condition C identical to Condition B for 83.9\% of the problems. Z3 counterexample feedback would be most valuable on problems where base pass@1 is well below the ceiling, leaving genuine headroom for guided repair.

\subsubsection{Limitations}

\textbf{L1: MBPP+ pass@1 comparison not completed.} All quantitative pass@1 claims (+8.94pp) are HumanEval+-specific. The h-m3 MBPP+ experiment was started (122/378 tasks, seed=42, Condition A only) but not completed due to compute constraints. Given the 0\% mypy error rate on MBPP+ (established by h-e1 and h-m1), Condition B is unlikely to show improvement on MBPP+, but this remains unconfirmed. All pass@1 claims in this paper are restricted to HumanEval+.

\textbf{L2: Type-specificity mechanism alternative unverified.} The general context enrichment interpretation is proposed as the most parsimonious explanation for the h-m3 result, but the ablation study needed to confirm it --- Condition B-Null: same prompt length as Condition B but with generic filler text instead of mypy output --- was not executed. Without this ablation, it is not possible to distinguish mypy content from mypy prompt length as the active ingredient. This is the highest-priority future experiment for mechanism confirmation.

\textbf{L3: Single model (GPT-4o-mini).} All experiments use GPT-4o-mini. The repair benefit may depend on the model's ability to interpret structured diagnostic text. Replication with at least one open-source code model and one stronger API model is needed to assess generalizability.

\textbf{L4: Z3 feedback tested only at high base pass@1.} The arithmetic subset experiment was conducted at base pass@1\_B = 86.6\%, leaving 13.4\% headroom. The null result for Z3 CE feedback may not generalize to settings with lower base pass@1, where CE signal would have more headroom to guide repair.

\subsubsection{Broader Implications}

This work establishes that adding mypy to LLM repair loops provides substantial, reproducible improvement in Python code generation quality on HumanEval+ (GPT-4o-mini, k=5). The practical cost is low: one subprocess call per repair round, approximately 30ms latency, no changes to model weights or architecture. The finding that the improvement is not type-specific --- if the general context enrichment interpretation holds --- suggests that other structured feedback channels deserve controlled experimental evaluation rather than being bundled with execution feedback and assumed beneficial. The Z3 pipeline (91.1\% valid specs) establishes feasibility for automated formal specification generation on algorithmic Python problems.

